%قبل از شروع مسئله ترجیحاً چیزی ننویسید. اما اگر خواستید بنویسید، در قالب لیستی (فقرات و شمارش) بنویسید.
%برای اضافه کردن نمره، بعدی از شروع‌مسئله، [x] را اضافه کنید و به جای x نمره مورد نظر را قرار دهید. برای زیرمسئله‌ها هم به همین صورت است.
%\شروع‌مسئله[10]
\شروع‌مسئله [12]
%\متن‌سیاه{سوال عملی اول:}\\
%\متن‌سیاه{دراپ‌تست}\\

%\متن‌سیاه{ورودی:}\\
%در خط اول ورودی، ابتدا k و سپس m داده می‌شود.\\
%\متن‌سیاه{خروجی:}\\
%در تنها سطر خروجی، تعداد آیفون‌های مورد نیاز می‌آید. اگر در هیچ صورتی چنین کاری ممکن نبود ۱- چاپ می‌شود.\\
%\متن‌سیاه{محدودیت‌ها:}\\
%$1\leq k\leq 10^4$

%$m\leq10^9$
%\\
%\textbf{\متن‌سیاه{ورودی و خروجی نمونه}\\
%% تست اول
%\begin{latin}
% \begin{tabularx}{\textwidth}{|X | X|}
% \hline
% \textbf{stdin} &
% \textbf{stdout}\\
% \hline
% 100 100&
% 1\\
% \hline
% 
% \end{tabularx}
%\end{latin}
%
%% تست دوم
%
%\begin{latin}
% \begin{tabularx}{\textwidth}{|X | X|}
% \hline
% \textbf{stdin} &
% \textbf{stdout}\\
% \hline
% 340 20&
% 3\\
% \hline
% 
% \end{tabularx}
%\end{latin}}
% صورت اصلی مسئله را در این‌جا بنویسید.
به پرسش‌های زیر درباره درخت‌های دودویی جست‌وجو پاسخ دهید.
%اگر مسئله چند قسمت داشت به این صورت عمل کنید:
\شروع{زیرمسائل}
\زیرمسئله[6] % مسئلۀ ۱آ
رابطه‌ای بازگشتی برای تعداد ددج‌های متفاوتی که می‌توان با اعداد
$a_1 < a_2 < ... < a_n$
نوشت به دست آورید. حال فرض کنید یک ددج ثابت داریم. رابطه‌ای بازگشتی برای تعداد دنباله‌های متفاوتی از اعضای این ددج را بیابید که درج آن‌ها منجر به ساخت این ددج شود.

\زیرمسئله[6] % مسئلۀ ۱ب
فرض کنید برای ساخت یک دنباله از اعضای یک ددج، به ترتیب از بیشترین عمق به کمترین عمق و در هر عمق از چپ به راست اعداد داخل گره‌های ددج را می‌نویسیم. آیا به ازای هر دنباله متناهی از اعداد صحیح متمایز، ددج متناظر با آن وجود دارد؟ درصورتی که مطمئن باشیم حداقل یک ددج متناظر با یک دنباله وجود دارد، آیا این ددج یکتاست؟

%\زیرمسئله % مسئلۀ ۱پ
%\زیرمسئله % مسئلۀ ۱ت
\پایان{زیرمسائل}

\solution{
	آ) برای هر درخت دودویی ریشه‌دار (که جهت چپ و راست آن نیز متمایز است) با
	$n$
	گره تنها یک ترتیب از مقداردهی 
	$a_1, ..., a_n$
	وجود دارد که درخت را تبدیل به ددج کند و هر ددج نیز یک درخت دودویی ریشه‌دار است بنابراین کافیست تعداد این درخت‌ها را بشماریم. اگر تعداد گره‌های زیردرخت چپ برابر با
	$i$
	باشد تعداد حالات ساخت درخت برابر است با
	$T(i)T(n - 1 - i)$
	بنابراین تعداد کل حالات برابر است با
	$$T(n) = \sum_{i=0}^{n-1}T(i)T(n-1-i)$$
	
	از آنجایی که عضو اول دنباله درج حتما ریشه است مکان آن در دنباله به طور یکتا مشخص می‌شود. اگر تعداد دنباله‌های درج را با
	$f$،
	تعداد دنباله‌های درج برای زیردرخت‌های چپ و راست را به ترتیب با
	$f_l$ و $f_r$
	و تعداد گره‌های زیردرخت‌های چپ و راست را به ترتیب با
	$n_l$ و $n_r$
	نمایش دهیم آنگاه برای ادامه دنباله درج باید ابتدا مکان $n_l$ گره زیردرخت چپ در دنباله را انتخاب کرده و سپس تعداد حالات مرتب کردن گره‌های زیردرخت چپ و زیردرخت راست را بشمریم بنابراین نهایتا خواهیم داشت
	$$f = \binom{n_l + n_r}{n_l}f_lf_r$$
	
	ب) به ازای هر دنباله دلخواه چنین ددج‌ای الزاما وجود ندارد. برای مثال دنباله 
	$3, 1, 2$
	را در نظر بگیرید. ۲ حتما باید ریشه باشد. ۱ نیز حتما فرزند ۲ است و چون از ۲ کمتر است فرزند چپ ۲ است. از آنجایی که فرزند چپ ریشه بلافاصله قبل از ریشه آمده پس ریشه فرزند راستی ندارد. در این صورت هیچ مکانی برای قرار دادن ۳ نداریم. بنابراین ددج‌ای متناظر با دنباله فوق وجود ندارد.
	
	در صورتی که از وجود یک ددج برای یک دنباله مطمئن باشیم، می‌توانیم ثابت کنیم چنین ددج‌ای یکتاست. برای اینکار ابتدا نشان می‌دهیم اگر یک ددج برای دنباله
	$n$
	تایی وجود داشته باشد، برای هر پسوند
	$k$
	تایی این دنباله نیز یک ددج وجود دارد. برای اینکار کافیست مرتبا راس‌هایی که در ابتدای لیست می‌آیند را حذف کنیم. حال با کمک استقرا یکتایی این ددج برای همه
	$k$
	ها از جمله
	$n$
	را نشان می‌دهیم. فرض کنید ددج متناظر با پسوند
	$k$تایی
	یکتا باشد. در این صورت اگر از هر ددج متناظر با پسوند
	$k+1$تایی
	گره
	$k+1$اُم
	را حذف کنیم به یک ددج یکتا می‌رسیم. سپس باید گره
	$k+1$اُم
	را به عنوان یک برگ در این ددج یکتا درج کنیم که حاصل یک ددج یکتا می‌شود. بنابراین ددج متناظر با پسوند 
	$k + 1$تایی
	نیز یکتاست. پایه استقرا نیز حالتی است که فقط ریشه وجود دارد که بدیهتا یکتاست. بنابراین با کمک استقرا ثابت کردیم اگر برای یک دنباله
	$n$تایی
	حداقل یک ددج معادل وجود داشته باشد، این ددج معادل یکتاست.

}




%%%%%%%%%%%%%%%%%%%%%%%%%%%%%%%%%%%%%%%%%%%%%%%%%%%%%
% برخی دستورات مفید برای تمرین‌ها. %
%%%%%%%%%%%%%%%%%%%%%%%%%%%%%%%%%%%%%%%%%%%%%%%%%%%%%

% فرمول‌های نمونه:
%$f(x) = \Theta(x)$
%$T(x, y) \leq \log(x) + 2^y + \binom{2n}{n}$
% {\tt code\_function}


% لیست موارد بدون شماره‌گذاری:
%\شروع{فقرات}
% \فقره مورد اول 
% \شروع{فقرات}
% \فقره مورد اول داخلی 
% \شروع{فقرات}
% \فقره اولین مورد خیلی داخلی 
% \فقره دومین مورد خیلی داخلی 
% \پایان{فقرات}
% \فقره مورد دوم داخلی 
% \پایان{فقرات}
% \فقره مورد دوم 
%\پایان{فقرات}

% لیست موارد شماره‌گذاری شده:
%\شروع{شمارش}
% \فقره یک مورد 
% \فقره یک مورد دیگر 
% \فقره باز هم یک مورد دیگر
%\پایان{شمارش}

% برای نوشتن یک فرمول چند خطی که یک شماره بخوره. برای حالت‌های دیگر به سایت زیر مراجعه کنید:
% https://www.overleaf.com/learn/latex/Aligning_equations_with_amsmath
%\begin{align} 
% \begin{split}
% (x+y)^3 &= (x+y)^2(x+y) \\
% &= (x^2+2xy+y^2)(x+y) \\
% &= (x^3+2x^2y+xy^2) + (x^2y+2xy^2+y^3) \\
% &= x^3+3x^2y+3xy^2+y^3
% \end{split} 
%\end{align}

% نوشتن ماتریس و آرایه دوبعدی
%\begin{align}
% A = 
% \begin{bmatrix}
% A_{11} & A_{21} \\
% A_{21} & A_{22}
% \end{bmatrix}
%\end{align}

% اضافه کردن تصویر با شرح. برای اضافه کردن تصویر بدون شرح، فقط دستور درج‌تصویر را استفاده کنید.
%\شروع{شکل}[h!]
%\تنظیم‌ازوسط
% \درج‌تصویر[پهنا=0.5\پهنای‌متن]{img.jpg}
%\شرح{شرح تصویر}
%\برچسب{تصویر:۱}
%\پایان{شکل}
%\پایان{زیرمسائل}

%کد زیر را ببینید.
%%\lstinputlisting{source_filename.py}
%\begin{latin}
%\begin{lstlisting}
%# Import the modules
%import sys
%import random
%
%ans = True
%
%while ans:
%question = raw_input("Ask the magic 8 ball a question: (press enter to quit) ")
%
%answers = random.randint(1,8)
%
%if question == "":
%sys.exit()
%
%elif answers == 1:
%print "It is certain"
%
%elif answers == 2:
%print "Outlook good"
%
%elif answers == 3:
%print "You may rely on it"
%
%elif answers == 4:
%print "Ask again later"
%
%elif answers == 5:
%print "Concentrate and ask again"
%
%elif answers == 6:
%print "Reply hazy, try again"
%
%elif answers == 7:
%print "My reply is no"
%
%elif answers == 8:
%print "My sources say no"
%\end{lstlisting}
%\end{latin}
